\statement{A.138}{Proved}{The entire transverse annular intersection and the literal handle-swap exclusion}
In the actual full profinite elliptic group \(G=\widehat B_{1,3}\),
put
\[
 H_X=\overline{\langle X,\sigma,\tau\rangle},\qquad
 H_Y=\overline{\langle Y,\sigma,\tau\rangle},\qquad
 H_B=\overline{\langle\sigma,\tau\rangle}.
\]
Then
\[
                         H_X\cap H_Y=H_B .
 \tag{A.138.1}
\]
Here \(H_B\) is the actual closed braid subgroup; no identification
with an abstract completion is asserted. Let \(\varepsilon\) be
the literal geometric handle swap
\[
 X\longleftrightarrow Y,\qquad
 \sigma\longmapsto\sigma^{-1},\qquad
 \tau\longmapsto\tau^{-1}.
 \tag{A.138.2}
\]
For every profinite commutator word \(w\) satisfying the ORIGINAL
marked four-strand pentagon, its A.121/A.131 cubic satisfies
\[
             \varepsilon(\Pi_w)\in H_X
                      \quad\Longleftrightarrow\quad w=1 .
 \tag{A.138.3}
\]
Thus this specified swap does not identify a nontrivial pentagon
word's negative cubic with the canonical positive word in \(H_X\).
No nontriviality of the actual canonical \(u\), or impossibility
of a correction outside \(H_X\), is asserted.
\begin{proof}
\emph{1. The actual based maps and the rank-two intersection.}
The automorphism (A.138.2) is the literal based action of
\citet[Proposition~2.5(1) and Remark~2.8]{A_Enriquez2014}.
Its generator rows square to the identity and extend to the
full completion. Only that action is used; the discrete
classification in part (2) is not transferred to profinite groups.
In particular \(H_Y=\varepsilon(H_X)\).

Write \(P=\ker(G\to S_3)\) and let
\(\pi_3:P\to E=\widehat F(X,Y)\) be the actual based deletion.
This map is used on \(P\), not on all of \(G\).
Put \(K=[Y,X]\), with commutator convention \(aba^{-1}b^{-1}\),
and define
\[
 A=\overline{\langle X,K\rangle},\qquad
 B_0=\overline{\langle Y,K\rangle}\quad\text{in }E.
 \tag{A.138.4}
\]
Both displayed pairs are actual free profinite bases. For example
\(\langle X,K\rangle=\langle X,YXY^{-1}\rangle\) is a free
factor of the index-two \(Y\)-parity subgroup with basis
\(X,YXY^{-1},Y^2\). The other pair is obtained from the
index-two \(X\)-parity subgroup. These are finite-cover
free-factor embeddings, as in A.134.

Let \(E_0\) be the kernel of the two exponent parities
\(E\to C_2^2\). The transversal \(\{1,X,Y,XY\}\) gives
the five non-tree Schreier edges
\[
 a_0=X^2,\quad b_0=Y^2,\quad k_0=K,\quad
 c_0=XYXY^{-1},\quad e_0=XY^2X^{-1}.
 \tag{A.138.5}
\]
Replace the last two by
\[
 j_0=c_0a_0^{-1}=XKX^{-1},\qquad
 l_0=b_0e_0^{-1}k_0^{-1}=YKY^{-1}.
 \tag{A.138.6}
\]
The inverse substitutions are
\(c_0=j_0a_0\), \(e_0=k_0^{-1}l_0^{-1}b_0\).
These are reversible Nielsen changes, so
\(a_0,b_0,k_0,j_0,l_0\) is a free profinite basis of \(E_0\).
Every element of \(A\cap B_0\) has both exponent parities zero.
The two index-two Schreier calculations give
\[
 A\cap E_0=\overline{\langle a_0,k_0,j_0\rangle},\qquad
 B_0\cap E_0=\overline{\langle b_0,k_0,l_0\rangle}.
 \tag{A.138.7}
\]
The continuous retraction of \(E_0\) fixing
\(a_0,k_0,j_0\) and killing \(b_0,l_0\) is the identity
on the first group and maps the second onto
\(\overline{\langle k_0\rangle}\). Hence
\[
 A\cap B_0=\overline{\langle K\rangle},\qquad
 A\cap B_0'=1.
 \tag{A.138.8}
\]
For the second equality, \(K\) is a primitive coordinate of
the abelianization in the actual basis \((Y,K)\) of \(B_0\);
therefore \(\overline{\langle K\rangle}\cap B_0'=1\).
This argument uses a continuous retraction in an explicit
finite-index group, not a general intersection theorem for
infinite-index completions.

\par\smallskip\noindent
\emph{2. The exact rank-three fiber and its swapped image.}
Set
\[
\begin{gathered}
 a=\sigma^2,\quad b=\tau^2,\quad m=\tau^{-1}a\tau,\qquad
 S_X=H_X\cap P,\quad S_Y=H_Y\cap P,\\
 h=X_3,\quad k=Y_3,\qquad
 N=\ker\pi_3=\widehat F(h,k,m),\qquad
 w_0=k^{-1}h^{-1}kh .
\end{gathered}
\tag{A.138.9}
\]
The actual fiber is that of A.128/A.134. Its corrected based
identity is
\[
                              bm=w_0 .
\tag{A.138.10}
\]
The completed source fiber/base diagram of A.134 is exact.
Its source fiber generators map to \((h,m,b)\), and its
source base maps injectively to the actual free pair \((X,K)\).
Consequently
\[
 N\cap S_X=\overline{\langle h,m,w_0\rangle},\qquad
 \pi_3(S_X)=A .
\tag{A.138.11}
\]
These are equalities for the actual compact groups. In
particular the kernel equality uses the already proved
completed exact sequence and base injection, not just a list
of discrete fiber generators.

On the six original pure generators \((X,Y,X_2,Y_2,a,b)\),
the map \(\pi_3\varepsilon\) has values
\[
                 (Y,X,Y^{-1},X^{-1},K^{-1},1).
\tag{A.138.12}
\]
They are exactly \(\varepsilon_2\pi_3\), where
\(\varepsilon_2(X)=Y\), \(\varepsilon_2(Y)=X\).
The six generators topologically generate \(P\), so
\(\pi_3\varepsilon=\varepsilon_2\pi_3\) on all of \(P\).
In particular \(\varepsilon(N)=N\).
The complete fiber rows are
\[
 \varepsilon(h)=k,\quad \varepsilon(k)=h,\quad
 \varepsilon(w_0)=w_0^{-1},\quad
 \varepsilon(m)=w_0m^{-1}w_0^{-1}.
\tag{A.138.13}
\]
Indeed \(\varepsilon(m)=\tau a^{-1}\tau^{-1}
=bm^{-1}b^{-1}\); substituting \(b=w_0m^{-1}\)
gives the last row. The handle rows follow from the actual
definitions of \(X_3,Y_3\). Therefore
\[
 N\cap S_Y=\overline{\langle k,m,w_0\rangle},\qquad
 \pi_3(S_Y)=B_0 .
\tag{A.138.14}
\]

\par\smallskip\noindent
\emph{3. The entire fiber, pure and nonpure intersections.}
Put \(x=h^{-1}\), \(y=k^{-1}\); then
\(N=\widehat F(x,y,m)\) and \(w_0=[y,x]\).
Let \(N_0\) be the index-four kernel of the \(x,y\)
exponent parities, with no condition on the \(m\) parity.
The transversal \(\{1,x,y,xy\}\) gives the nine basis elements
\[
\begin{gathered}
 a_1=x^2,\quad b_1=y^2,\quad w_0,\quad
 c_1=xyxy^{-1},\quad e_1=xy^2x^{-1},\\
 m,\quad m_x=xmx^{-1},\quad m_y=ymy^{-1},\quad
 m_{xy}=xymy^{-1}x^{-1}.
\end{gathered}
\tag{A.138.15}
\]
The five handle letters are the calculation of part 1; the
other four non-tree edges are the \(m\) loops at the four
vertices. Replace
\[
 j_1=c_1a_1^{-1}=xw_0x^{-1},\qquad
 l_1=b_1e_1^{-1}w_0^{-1}=yw_0y^{-1}.
\]
The inverses \(c_1=j_1a_1\),
\(e_1=w_0^{-1}l_1^{-1}b_1\) show that the resulting
nine letters remain a free profinite basis.

Let \(F_X=\overline{\langle x,m,w_0\rangle}\) and
\(F_Y=\overline{\langle y,m,w_0\rangle}\).
Their displayed rank-three bases are the source fiber bases
in (A.138.11), (A.138.14), with the reversible substitution
\(b=w_0m^{-1}\). Their intersection lies in \(N_0\), and
their parity kernels have bases
\[
\begin{aligned}
 F_X\cap N_0&=\overline{\langle a_1,w_0,j_1,m,m_x\rangle},\\
 F_Y\cap N_0&=\overline{\langle b_1,w_0,l_1,m,m_y\rangle}.
\end{aligned}
\tag{A.138.16}
\]
The continuous retraction fixing the five letters in the first
line and killing \(b_1,l_1,m_y,m_{xy}\) maps the second
group onto \(\overline{\langle m,w_0\rangle}\).
Thus, including the immediate reverse containment,
\[
 (N\cap S_X)\cap(N\cap S_Y)
      =\overline{\langle m,w_0\rangle}
      =\overline{\langle m,b\rangle}.
\tag{A.138.17}
\]

If \(g\in S_X\cap S_Y\), then (A.138.8), (A.138.11)
and (A.138.14) give \(\pi_3(g)=K^z\) for a unique
\(z\in\widehat{\mathbb Z}\).
The common element \(a=\sigma^2\) has \(\pi_3(a)=K\).
Hence \(a^{-z}g\) is in the intersection (A.138.17), and
\[
                    S_X\cap S_Y=\overline{\langle a,b,m\rangle}.
\tag{A.138.18}
\]
The reverse containment holds because all three displayed
generators lie in both sources. Finally, for
\(g\in H_X\cap H_Y\), choose a word in the common
\(\sigma,\tau\) with the same \(S_3\) permutation as \(g\).
Left multiplication by its inverse puts \(g\) in
(A.138.18), which is contained in \(H_B\).
This proves (A.138.1); the opposite inclusion is immediate.

\par\smallskip\noindent
\emph{4. The entire transporter and the exact swap exclusion.}
For the original pentagon word \(w\), A.131 gives
\[
\begin{gathered}
 \Pi_w=p_w\Pi_0p_w^{-1},\qquad
 \Pi_0=\tau^{-1}\sigma^{-1}Y,\\
 p_w=B_w^{-1}UZ,\quad
 B_w=w(a,b),\quad U=w(X,K),\quad Z=w(X_2,b).
\end{gathered}
\tag{A.138.19}
\]
The entire transporter belongs to \(P'\) and to \(S_X\).
Its swap has the literal expression
\[
\begin{gathered}
 R_+=\varepsilon(\Pi_0)=\tau\sigma X,\qquad
 p_\varepsilon=\varepsilon(p_w)
      =w(a^{-1},b^{-1})^{-1}
          w(Y,K^{-1})w(Y_2,b^{-1}),\\
 \varepsilon(\Pi_w)=p_\varepsilon R_+p_\varepsilon^{-1}.
\end{gathered}
\tag{A.138.20}
\]
In particular its order is exactly three, and it has the
same permutation and torsion label as \(R_+\).
The actual deletion values
\[
 \pi_3(X)=X,\quad\pi_3(Y)=Y,\quad
 \pi_3(X_2)=X^{-1},\quad\pi_3(Y_2)=Y^{-1},
 \quad\pi_3(a)=K,\quad\pi_3(b)=1
\]
give
\[
                    \pi_3(p_\varepsilon)=w(Y,K^{-1}).
\tag{A.138.21}
\]
The other two factors in (A.138.20) have a trivial argument.
No deletion map is applied to the nonpure cubic itself.

A.134 identifies \(H_X\) with the full source \(T\).
The source classification in A.135 has precisely the two
nonidentity order-three classes. The source image
\(R_X=X\sigma\tau\) is conjugate to \(R_+\) inside
\(H_X\), by the exact identity
\[
 R_+=(X^{-1}\tau)(X\sigma\tau)(X^{-1}\tau)^{-1}.
\tag{A.138.22}
\]
The actual character from A.124,
\[
 \chi_X:G\longrightarrow C_3,\qquad
 X\longmapsto1,\quad Y,\sigma,\tau\longmapsto0 ,
\]
distinguishes \(R_+\) from \(R_+^{-1}\).
Suppose \(\varepsilon(\Pi_w)\in H_X\).
Its torsion class and character therefore supply
\(g\in H_X\) with
\(\varepsilon(\Pi_w)=gR_+g^{-1}\).
The permutations of the two specified cubics agree, so the
permutation of \(g\) centralizes the three-cycle.
Right multiplication by a suitable power of \(R_+\)
leaves the conjugacy unchanged and makes \(g\) pure.
A.124 gives \(C_G(R_+)\cap P=1\). The two pure
transporters \(g,p_\varepsilon\) are therefore equal.
Thus \(p_\varepsilon\in S_X\); it also belongs to \(S_Y\),
since \(p_w\in S_X\).

By (A.138.18), \(p_\varepsilon\) lies in the actual closed
pure braid subgroup, and its third deletion is a power of
\(K\). But (A.138.21) lies in \(B_0'\), since \(w\) is a
commutator word. The primitive \(K\) coordinate in
(A.138.8) forces \(w(Y,K^{-1})=1\).
The pair \((Y,K^{-1})\) is an actual free profinite basis
of \(B_0\), so the original formal word \(w\) is one.
Conversely \(w=1\) makes the swapped cubic \(R_+\in H_X\).
This proves (A.138.3).

For the original canonical \(u\), A.37 supplies the required
pentagon. The positive word \(C_{+,G}\) of A.134 lies
in \(H_X\); if \(u\ne1\), it therefore cannot equal
\(\varepsilon(\Pi_u)\), or any of its conjugates by \(H_X\).
Moreover
\(\Pi_u^*=B_u\Pi_uB_u^{-1}\) from A.133 has
\(\varepsilon(B_u)=u(a^{-1},b^{-1})\in H_X\).
Consequently
\[
 \varepsilon(\Pi_u^*)\in H_X
       \quad\Longleftrightarrow\quad
 \varepsilon(\Pi_u)\in H_X
       \quad\Longleftrightarrow\quad u=1 .
\tag{A.138.23}
\]
No assertion that the actual \(u\) is nontrivial is required.
Corrections outside the specified annular group remain
unrestricted by this theorem.
\end{proof}

\paragraph{Source consolidation.}
The complete subgroup arguments consolidate
\path{work_A/actual_handle_swap_annular_exclusion38.md}
(\texttt{44b897d7}) and
\path{work_A/actual_transverse_annular_intersection38.md}
(\texttt{074eebb2}). Both actual finite-index free-basis
retractions and the source torsion normalization are retained.
The statement supplies no value of the canonical positive cube,
specified boundary or original CM character.

\paragraph{Unnumbered clarification: the specified boundary is the
original hexagon gate.}
The following complete ordered proof consolidates the accepted
note \path{work_A/actual_boundary_hexagon_equivalence38.md}
(\texttt{6770cf25}). Its decision content overlaps A.87 and
A.134; it is a proof clarification, not an additional vanishing
result or a separately numbered theorem.

Retain the ORIGINAL marking
\[
\begin{gathered}
 B=C_2(s)\amalg C_3(r),\qquad \delta=sr,\qquad
 F=\ker(B\to S_3)=\widehat F(x,y),\\
 x=\delta^{-2},\qquad y=sxs,\qquad z=(xy)^{-1},\qquad
 r(x,y,z)r^{-1}=(y,z,x).
\end{gathered}
\tag{A.138.C1}
\]
Let \(u\in F'\) be the actual canonical word, with proved
duality \(u(y,x)=u(x,y)^{-1}\), actual even \(k\), and
the actual full unit \(d=1+2k\) from A.103. Put
\[
\begin{gathered}
 u_0=u(x,y),\quad u_1=u(y,z),\quad u_2=u(z,x),\\
 v=u_0^{-1}s=su_0,\qquad D=v\delta^d,\qquad
 H_d=u_0x^k u_2z^k u_1y^k.
\end{gathered}
\tag{A.138.C2}
\]
Define the continuous substitution on the actual free group by
\[
 x'=\psi(x)=x^d,\qquad
 y'=\psi(y)=u_0^{-1}y^du_0,\qquad
 z'=\psi(z)=(x'y')^{-1}.
\tag{A.138.C3}
\]
No injectivity, surjectivity or third-peripheral property of
\(\psi\) is assumed. The SPECIFIED third-cusp condition is
\[
 {\cal C}:\qquad z'=x^k u_2z^d u_2^{-1}x^{-k}.
\tag{A.138.C4}
\]
We prove directly
\[
             {\cal C}\quad\Longleftrightarrow\quad
             D^3=1\quad\Longleftrightarrow\quad H_d=1 .
\tag{A.138.C5}
\]

\begin{proof}
\emph{Ordered modular identities.}
Since \(s\) swaps \(x,y\), duality gives
\(su_0s=u_0^{-1}\). Hence \(v^2=1\).
The relation \(\delta^2=x^{-1}\), the fixed parameter
\(d=1+2k\), and \(s\delta=r\) give
\[
 \delta^d=\delta x^{-k},\qquad
 D=u_0^{-1}r x^{-k}=u_0^{-1}y^{-k}r .
\tag{A.138.C6}
\]
These identities hold for integers and extend to the profinite
exponents by continuity. The \(S_3\) image of \(D\) is \(r\):
\(u_0\) is pure, \(\delta\) has order two there, and \(d\)
is odd. Thus \(D^3\in F\).

Expand the last word in (A.138.C6) three times with the
original cyclic action:
\[
\begin{aligned}
 D^3&=u_0^{-1}y^{-k}u_1^{-1}z^{-k}u_2^{-1}x^{-k},\\
 (D^3)^{-1}&=x^k u_2z^k u_1y^k u_0,\\
 H_d&=u_0(D^3)^{-1}u_0^{-1}.
\end{aligned}
\tag{A.138.C7}
\]
This is the exact ordered A.87 identity, including its inverse
and conjugation. It proves the last equivalence in
(A.138.C5) without an automorphy assumption.
The two required conjugations are equally literal:
\[
 Dx'D^{-1}=v x^dv^{-1}=u_0^{-1}y^du_0=y',
\tag{A.138.C8}
\]
because \(\delta^d\) commutes with \(x^d\), and
\[
\begin{aligned}
 D^{-1}x'D
 &=r^{-1}y^k u_0x^du_0^{-1}y^{-k}r\\
 &=x^k u_2z^du_2^{-1}x^{-k}.
\end{aligned}
\tag{A.138.C9}
\]
Here \(r^{-1}yr=x\), \(r^{-1}u_0r=u_2\), and
\(r^{-1}xr=z\). Thus \({\cal C}\) is precisely
\(Dz'D^{-1}=x'\), with its specified conjugator.

\par\smallskip\noindent
\emph{The cycle and the actual centralizers.}
Assume \({\cal C}\). Conjugating \(x'y'z'=1\) by \(D\)
and using (A.138.C8)--(A.138.C9) gives
\[
 y'(Dy'D^{-1})x'=1,\qquad
 Dy'D^{-1}=y'^{-1}x'^{-1}=z'.
\]
Thus \(D\) cycles \(x',y',z'\), and \(D^3\in F\)
centralizes both \(x',y'\).
Because \(d\) is a unit,
\(x=(x^d)^{d^{-1}}\), \(y=(y^d)^{d^{-1}}\).
The actual primitive centralizers of C.5 give
\[
 C_F(x^d)=\overline{\langle x\rangle},\qquad
 C_F(u_0^{-1}y^du_0)
      =u_0^{-1}\overline{\langle y\rangle}u_0.
\tag{A.138.C10}
\]
Their intersection is trivial: an equality
\(x^a=u_0^{-1}y^bu_0\) abelianizes in the actual free
coordinates to \((a,0)=(0,b)\), hence \(a=b=0\).
Both procyclic coordinates are faithful. Therefore \(D^3=1\).
This direction neither requires \(x',y'\) to generate \(F\)
nor cancels any endomorphism.

\par\smallskip\noindent
\emph{The reverse map exists only after the cube holds.}
Conversely assume \(D^3=1\). Together with \(v^2=1\),
the universal property of the full free profinite product
gives a genuine continuous homomorphism
\[
 \Psi:B\longrightarrow B,\qquad
 \Psi(s)=v,\quad\Psi(r)=D .
\]
It has
\[
 \Psi(\delta)=vD=\delta^d,\qquad
 \Psi(x)=x^d=x',\qquad
 \Psi(y)=vx'v^{-1}=y' .
\tag{A.138.C11}
\]
Thus its restriction to \(F\) is exactly the substitution
(A.138.C3). Apply this actual map to \(rzr^{-1}=x\).
It yields \(Dz'D^{-1}=x'\), which (A.138.C9) identifies
with the SPECIFIED equation \({\cal C}\).
No full automorphy of \(\Psi\) or \(\psi\) is used.
This completes (A.138.C5).

\par\smallskip\noindent
\emph{The original A.122 boundary and its full embedding.}
In \(E=\widehat F(X,Y)\), put
\[
 p=[Y,X],\quad q=X,\quad z_E=(pq)^{-1}=YX^{-1}Y^{-1}.
\]
The pair \((p,q)\) is an actual free profinite basis of
\(H=\overline{\langle p,q\rangle}\): it is obtained from
the free factor \(\langle X,YXY^{-1}\rangle\) of the
index-two \(Y\)-parity subgroup by replacing
\(YXY^{-1}\) with \(p=(YXY^{-1})X^{-1}\).
This is the A.122 finite-cover embedding, not an
infinite-index completion assumption.
Set
\[
\begin{gathered}
 f=u(p,q),\quad U=u(q,p)=f^{-1},\quad V=u(z_E,p),\\
 \xi=Uq^dU^{-1}=f^{-1}q^df,\qquad
 \omega=p^kVYf,\qquad
 \eta(X)=\xi,\quad\eta(Y)=\omega .
\end{gathered}
\tag{A.138.C12}
\]
Using \(YqY^{-1}=pq=z_E^{-1}\), calculate the full
ordered commutator:
\[
\begin{aligned}
 [\omega,\xi]
 &=p^kVYf(f^{-1}q^df)f^{-1}Y^{-1}V^{-1}p^{-k}\xi^{-1}\\
 &=p^kVz_E^{-d}V^{-1}p^{-k}\xi^{-1}.
\end{aligned}
\tag{A.138.C13}
\]
Consequently
\[
\begin{aligned}
 [\omega,\xi]=p^d
 &\Longleftrightarrow
 p^kVz_E^{-d}V^{-1}p^{-k}=p^d\xi\\
 &\Longleftrightarrow
 \psi(z_E)=p^kVz_E^dV^{-1}p^{-k},
\end{aligned}
\tag{A.138.C14}
\]
where \(\psi(p)=p^d\), \(\psi(q)=\xi\).
This is exactly (A.138.C4) under
\((x,y,z)=(p,q,z_E)\).
The free-basis identification is injective, so vanishing
of the evaluated \(H_d(p,q)\) is equivalent to vanishing
of the ORIGINAL \(H_d(x,y)\). We have proved
\[
                    [\omega,\xi]=p^d
                       \quad\Longleftrightarrow\quad H_d=1 .
\tag{A.138.C15}
\]

A.122 already proves
\(\Theta\eta=\eta c_{\rm ell}\) if and only if this
boundary equation holds, without cancelling \(\eta\).
A.123 gives full automorphy of \(\eta\) and its original
free-factor restriction conditionally on that boundary.
Finally A.134, at the original full-GT scope, supplies
the remaining equivalences
\[
\begin{aligned}
 \Theta\eta=\eta c_{\rm ell}
 &\Longleftrightarrow [\omega,\xi]=p^d
 \Longleftrightarrow H_d=1\\
 &\Longleftrightarrow C_{+,S}^{\,3}=1
 \Longleftrightarrow C_{+,G}^{\,3}=1 .
\end{aligned}
\tag{A.138.C16}
\]
The A.134 reverse direction uses the actual positive pair
\((d,u)\in\widehat{GT}\) only after the assumed hexagon,
C.30 and A.37 establish that membership. The primary
elliptic cube then lifts through the proved source embedding.
No prospective pair is treated as already acting.
\end{proof}

\paragraph{The original decision remains incomplete.}
A.87 already relates the original modular cube and cusp
obstruction. The clarification supplies the direct
cycle-and-centralizer proof for the exact A.122 marking,
avoiding a circular appeal to automorphy or descent.
Neither a negative elliptic cubic nor a coherent torsion
transporter replaces \(C_+\). The original pentagon has
not been proved to force any condition in (A.138.C16).
The canonical boundary, original hexagon, positive cube,
Main A and original CM decision remain Open.
